\section{Weekly Summary}

\begin{tcolorbox}[colback=yellow!10,colframe=black!50,title=AI-Generated Content]
\noindent The summary below was written by Claude (Anthropic) from that week's regression outputs and series descriptions. It has not been reviewed or fact-checked by a human, and should be read as an automated, informal recap -- not analysis.
\end{tcolorbox}

Welcome back to another week of the Random Number Casino, where we blindfold ourselves, throw darts at the FRED database, and pretend the resulting scatterplots mean something. As always, the house rules: these series are paired completely at random, no theory, no causation, no dignity. If two lines happen to hold hands, it's coincidence until proven otherwise -- and it is almost never proven otherwise.

We opened Sunday with the crown jewel of comedy: the World Pandemic Uncertainty Index for New Zealand versus U.S. private construction spending on communications. A truly inspired pairing. The result? Nothing. An R-squared of 0.008, a p-value that rolled its eyes at us (0.32), and a detrended version that somehow managed to be even more of a shrug. New Zealand's pandemic anxieties, it turns out, do not predict how many phone lines America pours concrete around. Devastating for exactly nobody.

Then came the actual event of the week. Monday paired two Consumer Price Index series -- Northeast core inflation against U.S. transportation-services CPI -- and the raw regression lit up like a slot machine, R-squared of 0.97 and a p-value with more zeros than my bank account. "A-ha," you say, "trend mirage, both are prices, both drift up, we've seen this scam before." And normally you'd be right. But here's the plot twist that genuinely surprised me: after detrending, it STILL held up. R-squared of 0.84 and a p-value that's basically zero. Two CPI measures moving together even after you strip out the shared upward march is one of the very few times this project stumbles into something that isn't pure nonsense -- probably because, shockingly, prices of different things tend to be pushed around by the same underlying inflation. A rare "actually kind of real" moment. Frame it.

The rest of the week returned us to our regularly scheduled mediocrity. Tuesday's "Other Goods and Services" CPI versus medical-services prices posted a gaudy 0.98 raw R-squared that then vaporized completely once detrended (p-value 0.97, a number so close to random it should file for unemployment) -- textbook shared-trend illusion. Wednesday's government training benefits versus a discontinued loan rate found nothing raw but a suspiciously "significant" negative blip after detrending, which I'm choosing to file under "if you flip enough coins, one lands on its edge." Thursday, Friday, and Saturday -- Great Lakes financial-services spending, Mid-Atlantic natural gas prices versus San Francisco tech, and margin rates on Western mortgages versus Israeli recession flags -- all delivered the barely-there, borderline, don't-you-dare-quote-me results this project runs on. So the final tally: one genuinely respectable relationship, a pile of trend mirages, and a few statistical hiccups I refuse to take seriously. Same time next week, same disclaimers, same absurd exercise.
